\section{Results \& Findings}
\subsection{Multi Criteria Analysis}

\subsubsection{Echo}\label{subsubsec:evaluation-echo}

\paragraph{Communication Protocols}
Golang has a built-in HTTP client that can perform REST calls. \cite{bib:golang-http-client}
For WebSockets, the \emph{gorilla/websocket} package provides a fast, well-tested WebSocket implementation for Go. \cite{bib:gorilla-websocket}

Since this criterion is not assigned any weight, Echo meets the requirements and passes the criterion.

\paragraph{Community}
The official Echo GitHub repository has over 32.200 stars \cite{bib:echo-github}, there is no dedicated Echo Discord server but the Golang Discord server has over 42.000 members \cite{bib:golang-discord}, and there are at least 6.500 results on GitHub for Echo projects \cite{bib:echo-github-search}.

The scores for this criterion are shown in \cref{tab:multi-criteria-analysis-echo-community}.

\begin{htable}[
		caption = {Echo community evaluation},
		label = {tab:multi-criteria-analysis-echo-community},
	]{
		colspec={lllX},
	}
	Sub-criterion & Weight & Function & Normalized score \\
	Number of stars & 5 & $\bar{x} = \frac{N}{N_{\max}}$ & $\bar{x} = \frac{32200}{32200} = 1$ \\
	Number of Discord users & 5 & $\bar{x} = \frac{N}{N_{\max}}$ & $\bar{x} = \frac{42000}{44000} = 0.95$ \\
	Number of examples & 10 & $\bar{x} = \frac{N}{N_{\max}}$ & $\bar{x} = \frac{6500}{508000} = 0.01$ \\
\end{htable}

\[
	S_{\text{community}} =
	\frac{
		5 \cdot 1 +
		5 \cdot 0.95 +
		10 \cdot 0.01
	}{
		5 + 5 + 10
	} = 0.49
\]

The normalized total score for this criterion is \emph{0.49}.

\paragraph{Developer Experience}
The proficiency ratings are as follows:
\begin{itemize}
	\item \textbf{Moustafa (0.60)} - Has built a REST API using Golang in his free time before, and feels comfortable with the language and its patterns.
	\item \textbf{William (0.05)} - Has not got any experience with Golang or Echo, but did review a REST API built with it, which gave him some idea of how the framework works.
\end{itemize}

Debugging Go applications can be done inside GoLand using JetBrains' built-in debugger. \cite{bib:goland-ide}
The documentation for Echo contains an installation guide and covers routing, middleware, and request handling, and the Go standard library has its own comprehensive documentation website as well. \cite{bib:echo-documentation}

The scores for this criterion are shown in \cref{tab:multi-criteria-analysis-echo-developer-experience}.

\begin{htable}[
		caption = {Echo developer experience evaluation},
		label = {tab:multi-criteria-analysis-echo-developer-experience},
	]{
		colspec={lllX},
	}
	Sub-criterion & Weight & Function & Normalized score \\
	Learning curve & 20 & $\bar{x} = \frac{x_1 + x_2}{2}$ & $\bar{x} = \frac{0.60 + 0.05}{2} = 0.33$ \\
	Debugging support & 10 & Binary (0 or 1) & 1 \\
	Documentation quality & 15 & Binary (0 or 1) & 1 \\
\end{htable}

\[
	S_{\text{developer experience}} =
	\frac{
		20 \cdot 0.33 +
		10 \cdot 1 +
		15 \cdot 1
	}{
		20 + 10 + 15
	} = 0.70
\]

The normalized total score for this criterion is \emph{0.70}.

\paragraph{Ecosystem}
JetBrains provides GoLand, a dedicated IDE for the Go programming language, which supports plugins for linting, testing, and other development tools. \cite{bib:goland-ide}
The Go package repository at \texttt{pkg.go.dev} hosts a large number of packages compatible with Echo applications.

The scores for this criterion are shown in \cref{tab:multi-criteria-analysis-echo-ecosystem}.

\begin{htable}[
		caption = {Echo ecosystem evaluation},
		label = {tab:multi-criteria-analysis-echo-ecosystem},
	]{
		colspec={lllX},
	}
	Sub-criterion & Weight & Function & Normalized score \\
	IDE plugin support & 25 & Binary (0 or 1) & 1 \\
	Number of libraries & 15 & $\bar{x} = \frac{N}{N_{\max}}$ & $\bar{x} = \frac{0}{2000} = 0$ \\
\end{htable}

\[
	S_{\text{ecosystem}} =
	\frac{
		25 \cdot 1 +
		15 \cdot 0
	}{
		25 + 15
	} = 0.63
\]

The normalized total score for this criterion is \emph{0.63}.

\paragraph{Quality Assurance}
Echo makes use of the Golang testing platform, which is built into the standard library and supports unit and integration testing. \cite{bib:golang-testing}
The Go ecosystem provides linting tools such as \emph{golangci-lint}, which bundles multiple linters and integrates with major JetBrains IDEs. \cite{bib:golangci-lint}
The last minor release of the Echo framework was the last six months. \cite{bib:echo-github}

The scores for this criterion are shown in \cref{tab:multi-criteria-analysis-echo-quality-assurance}.

\begin{htable}[
		caption = {Echo quality assurance evaluation},
		label = {tab:multi-criteria-analysis-echo-quality-assurance},
	]{
		colspec={lllX},
	}
	Sub-criterion & Weight & Function & Normalized score \\
	Testing & 15 & Binary (0 or 1) & 1 \\
	Static code analysis & 5 & Binary (0 or 1) & 1 \\
	Regular updates & 5 & Binary (0 or 1) & 1 \\
\end{htable}

\[
	S_{\text{quality assurance}} =
	\frac{
		15 \cdot 1 +
		5 \cdot 1 +
		5 \cdot 1
	}{
		15 + 5 + 5
	} = 1.00
\]

The normalized total score for this criterion is \emph{1.00}.

\subsubsection{Spring Boot}\label{subsubsec:evaluation-spring-boot}

\paragraph{Communication Protocols}
Spring Boot provides a built-in \texttt{RestTemplate} and the newer \texttt{WebClient} for performing RESTful API calls. \cite{bib:spring-boot-rest}
For WebSocket connections, Spring Boot includes the \emph{spring-websocket} module, which supports full-duplex communication using the STOMP messaging protocol. \cite{bib:spring-boot-websocket}

Since this criterion is not assigned any weight, Spring Boot meets the requirements and passes the criterion.

\paragraph{Community}
The official Spring Boot GitHub repository has over 80.000 stars \cite{bib:spring-boot-github}. Although there is no official Spring Boot Discord server, there are other communities focused on Java in general. Additionally, a search on GitHub returns over 400,000 Spring Boot projects \cite{bib:spring-boot-github-search}.

The scores for this criterion are shown in \cref{tab:multi-criteria-analysis-spring-boot-community}.

\begin{htable}[
		caption = {Spring Boot community evaluation},
		label = {tab:multi-criteria-analysis-spring-boot-community},
	]{
		colspec={lllX},
	}
	Sub-criterion & Weight & Function & Normalized score \\
	Number of stars & 5 & $\bar{x} = \frac{N}{N_{\max}}$ & $\bar{x} = \frac{80000}{80000} = 1$ \\
	Number of Discord users & 5 & $\bar{x} = \frac{N}{N_{\max}}$ & $\bar{x} = \frac{0}{0} = 0.0$ \\
	Number of examples & 10 & $\bar{x} = \frac{N}{N_{\max}}$ & $\bar{x} = \frac{400000}{508000} = 0.79$ \\
\end{htable}

\[
	S_{\text{community}} =
	\frac{
		5 \cdot 1 +
		5 \cdot 0 +
		10 \cdot 0.79
	}{
		5 + 5 + 10
	} = 0.65
\]

The normalized total score for this criterion is \emph{0.65}.

\paragraph{Developer Experience}
The proficiency ratings are as follows:
\begin{itemize}
	\item \textbf{Moustafa (0.25)} - Has some experience with Java from school, but has not used Spring Boot or the Spring ecosystem before.
	\item \textbf{William (0.25)} - Has limited experience with Java and has not used Spring Boot, though the general structure is somewhat familiar from other frameworks.
\end{itemize}

Debugging Spring Boot applications can be done inside IntelliJ IDEA, which provides first-class support for the Spring framework including run configurations and bean inspection. \cite{bib:intellij-spring}
The Spring Boot documentation is extensive and covers installation, configuration, routing, security, and data access. \cite{bib:spring-boot-documentation}

The scores for this criterion are shown in \cref{tab:multi-criteria-analysis-spring-boot-developer-experience}.

\begin{htable}[
		caption = {Spring Boot developer experience evaluation},
		label = {tab:multi-criteria-analysis-spring-boot-developer-experience},
	]{
		colspec={lllX},
	}
	Sub-criterion & Weight & Function & Normalized score \\
	Learning curve & 20 & $\bar{x} = \frac{x_1 + x_2}{2}$ & $\bar{x} = \frac{0.25 + 0.25}{2} = 0.25$ \\
	Debugging support & 10 & Binary (0 or 1) & 1 \\
	Documentation quality & 15 & Binary (0 or 1) & 1 \\
\end{htable}

\[
	S_{\text{developer experience}} =
	\frac{
		20 \cdot 0.25 +
		10 \cdot 1 +
		15 \cdot 1
	}{
		20 + 10 + 15
	} = 0.67
\]

The normalized total score for this criterion is \emph{0.67}.

\paragraph{Ecosystem}
IntelliJ IDEA Ultimate provides dedicated Spring Boot support including Spring-specific tooling, run configurations, and bean diagrams. \cite{bib:intellij-spring}
The Maven Central registry contains over 686.000 packages, a large portion of which are compatible with Spring Boot applications. \cite{bib:maven-central}

The scores for this criterion are shown in \cref{tab:multi-criteria-analysis-spring-boot-ecosystem}.

\begin{htable}[
		caption = {Spring Boot ecosystem evaluation},
		label = {tab:multi-criteria-analysis-spring-boot-ecosystem},
	]{
		colspec={lllX},
	}
	Sub-criterion & Weight & Function & Normalized score \\
	IDE plugin support & 25 & Binary (0 or 1) & 1 \\
	Number of libraries & 15 & $\bar{x} = \frac{N}{N_{\max}}$ & $\bar{x} = \frac{686}{2000} = 0.34$ \\
\end{htable}

\[
	S_{\text{ecosystem}} =
	\frac{
		25 \cdot 1 +
		15 \cdot 0.34
	}{
		25 + 15
	} = 0.75
\]

The normalized total score for this criterion is \emph{0.75}.

\paragraph{Quality Assurance}
Spring Boot integrates with JUnit and provides the \texttt{@SpringBootTest} annotation for integration testing, as well as MockMvc for testing web layers in isolation. \cite{bib:spring-boot-testing}
The Java ecosystem offers mature static analysis tools such as Checkstyle, PMD, and SpotBugs, all of which integrate well with JetBrains IntelliJ IDEA. \cite{bib:java-linting}
The last minor release of Spring Boot was within the last six months. \cite{bib:spring-boot-github}

The scores for this criterion are shown in \cref{tab:multi-criteria-analysis-spring-boot-quality-assurance}.

\begin{htable}[
		caption = {Spring Boot quality assurance evaluation},
		label = {tab:multi-criteria-analysis-spring-boot-quality-assurance},
	]{
		colspec={lllX},
	}
	Sub-criterion & Weight & Function & Normalized score \\
	Testing & 15 & Binary (0 or 1) & 1 \\
	Static code analysis & 5 & Binary (0 or 1) & 1 \\
	Regular updates & 5 & Binary (0 or 1) & 1 \\
\end{htable}

\[
	S_{\text{quality assurance}} =
	\frac{
		15 \cdot 1 +
		5 \cdot 1 +
		5 \cdot 1
	}{
		15 + 5 + 5
	} = 1.00
\]

The normalized total score for this criterion is \emph{1.00}.

\subsubsection{NestJS}\label{subsubsec:evaluation-nestjs}

\paragraph{Communication Protocols}
NestJS runs on Node.js, so either the Fetch API or \emph{Axios} can be used to perform RESTful API calls. \cite{bib:axios}
For WebSocket connections, NestJS has a built-in WebSockets module that supports both native WebSockets and Socket.IO. \cite{bib:nestjs-websockets}

Since this criterion is not assigned any weight, NestJS meets the requirements and passes the criterion.

\paragraph{Community}
The official NestJS GitHub repository has over 75.000 stars \cite{bib:nestjs-github}, the official NestJS Discord server has over 51.000 members \cite{bib:nestjs-discord}, and there are at least 80.000 results on GitHub for NestJS projects \cite{bib:nestjs-github-search}.

The scores for this criterion are shown in \cref{tab:multi-criteria-analysis-nestjs-community}.

\begin{htable}[
		caption = {NestJS community evaluation},
		label = {tab:multi-criteria-analysis-nestjs-community},
	]{
		colspec={lllX},
	}
	Sub-criterion & Weight & Function & Normalized score \\
	Number of stars & 5 & $\bar{x} = \frac{N}{N_{\max}}$ & $\bar{x} = \frac{75000}{75000} = 1$ \\
	Number of Discord users & 5 & $\bar{x} = \frac{N}{N_{\max}}$ & $\bar{x} = \frac{51000}{51000} = 1$ \\
	Number of examples & 10 & $\bar{x} = \frac{N}{N_{\max}}$ & $\bar{x} = \frac{80000}{508000} = 0.16$ \\
\end{htable}

\[
	S_{\text{community}} =
	\frac{
		5 \cdot 1 +
		5 \cdot 1 +
		10 \cdot 0.16
	}{
		5 + 5 + 10
	} = 0.58
\]

The normalized total score for this criterion is \emph{0.58}.

\paragraph{Developer Experience}
The proficiency ratings are as follows:
\begin{itemize}
	\item \textbf{Moustafa (0.25)} - Has minimal experience with JavaScript frameworks such as VueJS and Svelte, but has not used NestJS or TypeScript extensively.
	\item \textbf{William (0.50)} - Has experience with several JavaScript and TypeScript frameworks and feels that the Angular-inspired structure of NestJS is relatively familiar.
\end{itemize}

Debugging NestJS applications can be done inside WebStorm, which provides full TypeScript and Node.js debugging support. \cite{bib:webstorm-ide}
The NestJS documentation contains an installation guide and covers modules, controllers, providers, middleware, and WebSockets in detail. \cite{bib:nestjs-documentation}

The scores for this criterion are shown in \cref{tab:multi-criteria-analysis-nestjs-developer-experience}.

\begin{htable}[
		caption = {NestJS developer experience evaluation},
		label = {tab:multi-criteria-analysis-nestjs-developer-experience},
	]{
		colspec={lllX},
	}
	Sub-criterion & Weight & Function & Normalized score \\
	Learning curve & 20 & $\bar{x} = \frac{x_1 + x_2}{2}$ & $\bar{x} = \frac{0.25 + 0.50}{2} = 0.38$ \\
	Debugging support & 10 & Binary (0 or 1) & 1 \\
	Documentation quality & 15 & Binary (0 or 1) & 1 \\
\end{htable}

\[
	S_{\text{developer experience}} =
	\frac{
		20 \cdot 0.38 +
		10 \cdot 1 +
		15 \cdot 1
	}{
		20 + 10 + 15
	} = 0.72
\]

The normalized total score for this criterion is \emph{0.72}.

\paragraph{Ecosystem}
WebStorm provides full support for TypeScript and Node.js development, and the JetBrains plugin marketplace includes NestJS-specific plugins for code generation and navigation. \cite{bib:webstorm-ide}
The npm registry contains over 2.000.000 packages, a large number of which are compatible with NestJS applications. \cite{bib:npm-registry}

The scores for this criterion are shown in \cref{tab:multi-criteria-analysis-nestjs-ecosystem}.

\begin{htable}[
		caption = {NestJS ecosystem evaluation},
		label = {tab:multi-criteria-analysis-nestjs-ecosystem},
	]{
		colspec={lllX},
	}
	Sub-criterion & Weight & Function & Normalized score \\
	IDE plugin support & 25 & Binary (0 or 1) & 1 \\
	Number of libraries & 15 & $\bar{x} = \frac{N}{N_{\max}}$ & $\bar{x} = \frac{2000}{2000} = 1$ \\
\end{htable}

\[
	S_{\text{ecosystem}} =
	\frac{
		25 \cdot 1 +
		15 \cdot 1
	}{
		25 + 15
	} = 1.00
\]

The normalized total score for this criterion is \emph{1.00}.

\paragraph{Quality Assurance}
NestJS integrates with Jest out of the box and provides utilities for unit testing individual components and end-to-end testing the full application. \cite{bib:nestjs-testing}
The TypeScript ecosystem provides ESLint with TypeScript-specific rules, and NestJS projects are scaffolded with ESLint configuration by default. \cite{bib:nestjs-linting}
The last minor release of NestJS was within the last six months. \cite{bib:nestjs-github}

The scores for this criterion are shown in \cref{tab:multi-criteria-analysis-nestjs-quality-assurance}.

\begin{htable}[
		caption = {NestJS quality assurance evaluation},
		label = {tab:multi-criteria-analysis-nestjs-quality-assurance},
	]{
		colspec={lllX},
	}
	Sub-criterion & Weight & Function & Normalized score \\
	Testing & 15 & Binary (0 or 1) & 1 \\
	Static code analysis & 5 & Binary (0 or 1) & 1 \\
	Regular updates & 5 & Binary (0 or 1) & 1 \\
\end{htable}

\[
	S_{\text{quality assurance}} =
	\frac{
		15 \cdot 1 +
		5 \cdot 1 +
		5 \cdot 1
	}{
		15 + 5 + 5
	} = 1.00
\]

The normalized total score for this criterion is \emph{1.00}.

\subsubsection{Overview}
\Cref{tab:score-results} shows an overview of the framework scores in comparison to one another.

\begin{htable}[
    caption = {Score results},
    label = {tab:score-results},
] {
    colspec = {Xllll}
}
    Criterion & Weight & Echo & Spring Boot & NestJS \\
    Communication Protocols & N/A & pass & pass & pass \\
    Community & 20 & $20 \cdot 0.49 = 9.80$ & $20 \cdot 0.65 = 13.00$ & $20 \cdot 0.58 = 11.60$ \\
    Developer Experience & 45 & $45 \cdot 0.70 = 31.50$ & $45 \cdot 0.67 = 30.15$ & $45 \cdot 0.72 = 32.40$ \\
    Ecosystem & 40 & $40 \cdot 0.63 = 25.20$ & $40 \cdot 0.75 = 30.00$ & $40 \cdot 1.00 = 40.00$ \\
    Quality Assurance & 25 & $25 \cdot 1.00 = 25.00$ & $25 \cdot 1.00 = 25.00$ & $25 \cdot 1.00 = 25.00$ \\
\end{htable}

The final ratings for each framework, normalized to a rating between 0 and 10, are presented in \cref{tab:scores}.

\begin{vtable}[
    caption = {Final scores},
    label = {tab:scores},
] {
    colspec = {XXXX}
}
    Framework & Echo & Spring Boot & NestJS \\
    Score & 7.04 & 7.55 & 8.38 \\
\end{vtable}

\subsection{Prototyping}

\subsubsection{Spring Boot}
Setting up and creating the project was relatively easy.
For building the prototype we used IntelliJ IDEA, which integrates well with Java and Spring Boot.
We decided to use Spring Boot 4 together with Java version 21.

Creating controllers, services, and repositories, as well as setting up the ORM, was straightforward due to the annotation-based approach provided by Spring Boot.
Dependency injection is handled out of the box, which simplifies the structure of the application.
Adding tests was also relatively easy because of the built-in support within the framework.

In addition, observability was added using Grafana, Prometheus, and Loki. Setting this up was mostly straightforward, as it mainly required adding the necessary dependencies and configuring the application.
There were some difficulties when configuring the security settings in Spring Boot, but after some experimentation the configuration became clearer.

The documentation for Spring Boot was generally sufficient.
Because Spring Boot is a mature and widely used framework, there is a large amount of information available online.
However, some resources were outdated, especially since this prototype uses the newer Spring Boot 4 version, which means certain solutions were deprecated.
Despite this, the official Spring Boot documentation provided enough examples and it was generally possible to find solutions to the problems encountered.

\subsubsection{NestJS}
Setting up the project was straightforward using the NestJS CLI.
For development we used a Docker setup with multiple containers (for monitoring, database, s3).
Getting hot reload working inside Docker took some extra effort, and in the end we could not get it working on Windows.
This was however, more of a DevOps issue than something related to NestJS itself.

Controllers, services and repositories were easy to set up using decorators.
Dependency injection is built into NestJS and worked well across the different layers.
This contributes to the developer experience, as it makes the structure predictable and easier to work with.

For persistence we used TypeORM with a repository pattern.
The developer environment was easy to set up with synchronization enabled.
NestJS strongly advised not to use synchronization in production but instead implement migrations.
This was doable, but proved to be a bit difficult.

The API was implemented as a NestJS controller and validation was done using Zod and NestJS pipes.
Validation caused some issues at the start.
The default approach in NestJS validates all controller parameters, while in our case only the request body needed validation.
Because of this, Zod object validation would fail when the endpoint also included route parameters.
This was resolved by adjusting the validation setup after looking up similar issues online.

One other issue we ran into was mapping database models to domain entities when relations were involved.
Including related entities could cause infinite recursion.
This was solved in a simple way by not including certain nested entities, which worked but was not ideal since it felt a bit `hacky'.

The NestJS documentation was clear and helpful.
It includes introductory chapters that made it easy to get started and build a solid foundation.
For more advanced topics, there are separate sections that go into more detail, which helped when implementing more complex parts of the application.
The documentation also highlights best practices, which can be easy to overlook but are useful as guidance during development.

We did not really use (or rather need) much of the community, besides a couple stackoverflow/medium posts, so it is hard to comment on that.
